\begin{rulebox}{Purpose}
This optional rule represents the loose exchange of fire by commanded bodies of
musketeers before the main bodies became fully engaged. It adds limited probing
and reaction fire without changing the normal combat system. A skirmish may
create an opportunity, but another unit must be positioned to exploit it before
the dice are rolled.
\end{rulebox}

\begin{multicols}{2}
\section{Optional Rule}

\subsection{1. Eligibility and timing}

At the beginning of a friendly March Phase, each friendly leader may designate
one ordered Foot unit of the same force (color) in its own or an adjacent hex.
A Foot unit may not be designated twice. Place a Skirmish---Return marker on
each designated unit. Only designated units may skirmish; they need not remain
near the leader after beginning movement. For example, Rupert and Byron could
designate one unit each (two total).

The target must be an ordered enemy unit that is not in a close, train, or woods
hex. Leaders, artillery pieces, and trains may neither initiate nor be targeted.

Skirmishing is permitted only during Game Turns 1--7 of the full game. It is not
used in the Short Game. No skirmishes may be declared during a flash-rain turn
\ruleref{5.4}.

An enemy unit may be selected as the target of no more than one skirmish in a
March Phase.

\subsection{2. Declaring a skirmish}

During its normal movement, an eligible Foot unit may enter a hex adjacent to an
eligible target and declare a skirmish. The unit pays the normal cost to enter
the adjacent hex and must then have at least \textbf{2 MPs remaining}: 1 MP to
skirmish and 1 MP to withdraw.

The player identifies the target and moves the designated unit's
\textbf{Skirmish---Return} marker into the hex just vacated; no other piece may
enter that hex. The skirmishing unit's movement ends. Complete all other
ordinary movement before rolling any skirmish dice. Remove unused designation
markers at the end of the March Phase.

\subsection{3. Resolving skirmishes}

After all ordinary movement in the March Phase is complete, resolve declared
skirmishes one at a time in any order chosen by the active player:

\begin{enumerate}
  \item The skirmishing unit fires at its declared target, and the target makes
        a \textbf{Reaction Roll}. Roll one die for each unit and consult the
        \textbf{6+} column of the Artillery Table \ruleref{17.2}. Other
        artillery rules do not apply: adjacency establishes eligibility, and no
        line of sight or range is traced.
  \item Record both results, but do not apply either result yet. The two attacks
        are simultaneous; a result suffered by one unit never cancels its own
        fire.
  \item Return the skirmishing Foot unit to the marked hex. This withdrawal
        costs the reserved 1 MP regardless of terrain.
  \item Apply both recorded results simultaneously, then remove the marker. A
        \textbf{Dd} result disrupts an ordered target; a dash has no effect.
\end{enumerate}

A Reaction Roll is not always literal return fire. For Foot it represents return
musketry; for Horse it may represent pistol or carbine fire, or a short
countercharge. Horse may be a target, but may not initiate a skirmish.

\subsection{4. After the exchange}

The March Phase ends after all marked skirmishes have been resolved. No unit may
move in response to a skirmish result. Proceed to the active player's Combat
Phase normally. Thus, a unit already positioned to attack may exploit a newly
disrupted enemy, but the active player cannot first see the result and then move
up an attacker.

Temporary adjacency created by the skirmishing unit has no effect after it
withdraws. All mandatory-attack requirements are determined from the units'
positions at the beginning of the ensuing Combat Phase \ruleref{9.0}.

\begin{rulebox}{Skirmish Fire: 6+ Column}
\centering
\begin{tabularx}{\linewidth}{@{}>{\bfseries}c Y@{}}
  \toprule
  Die Roll & Result \\
  \midrule
  1 & Dd---target disrupted \\
  2--6 & \dash---no effect \\
  \bottomrule
\end{tabularx}
\end{rulebox}

\begin{rulebox}{Rules References}
Movement and MP limits: \ruleref{8.1--8.8}. Artillery die-roll procedure and
table: \ruleref{7.0}, \ruleref{17.2}. Disruption: \ruleref{11.8}. Short Game:
\ruleref{16.0}.
\end{rulebox}

\section{Design Intent}

The idea here is to show some level of skirmishing at this scale while
preventing unrealistic combined arms attacks, so any Horse waiting to exploit a
disruption must be committed first. Most exchanges will do nothing, but
occasionally one will open a useful opportunity. Think of this as an amalgam of
skirmishing, advancing fire and retreating fire from something like GMT's
Musket \& Pike series.
\end{multicols}
